% !TeX root = ../教材题库_全部条目_转写工作稿.tex
% Source: XJ-B1-C04-EX-P139-Y0872
% 第4章 / 4.5 函数模型及其应用 / 教材 p.139、140 / PDF 第146、147页
% OCR未校对：公式与文字均可能有误；未完成数学TeX转写；禁止直接发给学生
% 原题及相邻原解截图（核对用）：../题目原图/XJ-B1-C04-EX-P139-Y0872-01.webp, ../题目原图/XJ-B1-C04-EX-P139-Y0872-02.webp
% 本文件只在转写完成后作为题干使用；OCR 条目尚不是可靠题干。
\begin{bankitem}{4.5 例题 例题1}{第4章 / 4.5 函数模型及其应用 / 教材 p.139、140 / PDF 第146、147页}{XJ-B1-C04-EX-P139-Y0872}{OCR未校对：公式与文字均可能有误；未完成数学TeX转写；禁止直接发给学生}
\begin{ocrdraft}
[教材p.139 / PDF第146页]
\par
0〔〔]比较3，一在和3一在区间匚0，十。。）上增长的快慢．
\par
解两函数的图象如图4，5一1所示．
\par
4
\par
0
\par
4
\par
8
\par
12
\par
图4．5一1
\par
开始3，一、0一路领先，但越来越慢；3，一一匀速前进；在一16处两者相等．
\par
当工>16时，>16两边开方得x>40，即一直有莖>、在，而且前者与后
\par
者之比越来越大，所以3一增长得更快．
\par
开运动会常常要分组选拔，函数赛跑也可以先分组比一比．我们考虑五类在区
\par
间（0，十亠）上单调递增的函数：
\par
指数函数3，一a'>1)算，A组；幂函数3，一>1)算B组；一次函数3-十
\par
>0）算C组；幂函数3，一的（0<仅<1）算D组；对数函数3，一log“（>1）算最后
\par
[教材p.140 / PDF第147页]
\par
一组，E组．
\par
同一组的比赛容易分出高低，看图便知分晓．
\par
从图4．5一2（D看出，A组内，a越大跑得越快；E组内，磊越小跑得越快．
\par
从图4．5一2（2）看出，B组和D组一起比赛，都是a越大跑得越快．
\par
厂3`）彐og凵攴
\par
y=xly=xl\ocrglyph{00E4}
\par
y=loggx
\par
（2）
\par
图4．5一2
\par
现在来看C组，一次函数3，一十b（>0）．
\par
如果两个一次函数的一次项系数相等，只有常数项不同，则两个函数的差是常
\par
数．起跑时在前面的永远在前面，领先距离永远不变．从图象上看，是两条平行
\par
直线．
\par
如果两个一次函数的一次项系数不相等，系数大的跑得就快．不管起跑时落后
\par
有多少，系数大的总能后来居上，而且将遥遥领先．在方格纸上画几个一次函数的
\par
图象便能看出这个规律．
\par
小组选拔赛的情形一目了然．组与组之间的比赛呢？
\par
上面已经对B，D两组做了比较．
\par
数学上可以证明，B组的任一个成员3，一（a>1)总比C组的3，一十氕差>0）
\par
增长得快．
\end{ocrdraft}
\end{bankitem}
